\section{Berechnung der Beiträge bei Tarifänderungen}\label{sec:tarifaenderung}

Dieser Abschnitt regelt die Beitragsberechnung bei individuellen
Tarifänderungen, bei Beitragsanpassungen für den Bestand sowie die
Beitragsermäßigungen im Alter.

\subsection{Beiträge bei individuellen Tarifänderungen}\label{sec:umstufung}

Bei individuellen Tarifänderungen (Tarifwechsel nach \S\,204 VVG,
Höher-/Tieferstufungen innerhalb des Tarifs, z.\,B. Änderung des
Selbstbehalts) wird der aus dem bisherigen Vertragsteil kalkulatorisch
zugeordnete Anteil der Alterungsrückstellung gemäß \S\,13 Abs.\,1 KVAV
vollständig prämienmindernd angerechnet. Die Berechnung folgt den
Formeln des Abschnitts~B der Anlage~1 zur KVAV (\S\S\,12, 13 KVAV) in
folgender Darstellung: Mit der Hilfsfunktion
\begin{equation}
b(u) = \frac{\bigl[12(1-\sigma-\Omega^{\mathrm{n}}_x)\,a^{\mathrm{n}}_x - \alpha^{\mathrm{n}}_x\bigr]\, b^{\mathrm{n}}_x
      \;-\; \bigl[12(1-\sigma-\Omega^{\mathrm{a}}_x)\,a^{\mathrm{a}}_x - \alpha^{\mathrm{a}}_x\bigr]\, b^{\mathrm{a}}_x
      \;+\; \bigl[12(1-\sigma-\Omega^{\mathrm{a}}_x)\,a^{\mathrm{a}}_x - u\bigr]\, b^{\mathrm{a}}}
      {12(1-\sigma-\Omega^{\mathrm{n}}_x)\,a^{\mathrm{n}}_x - u}
\label{eq:umstufung}
\end{equation}
ergibt sich der neue monatliche Beitrag als
\begin{equation}
b^{\mathrm{n}} = \max\{b(0);\, b(\tilde{\alpha})\},
\qquad
\hat{b}^{\mathrm{n}} = \max\bigl\{\min\{b^{\mathrm{n}};\, b^{\mathrm{n}}_x\};\, b^{\mathrm{n}}_{x_e}\bigr\},
\label{eq:umstufung2}
\end{equation}
wobei
\begin{itemize}\itemsep1pt
  \item die Kennzeichnungen $\mathrm{a}$/$\mathrm{n}$ den alten bzw. neuen Vertragsteil
        bezeichnen (bei Umstufungen innerhalb des Tarifs sind die
        Rechnungsgrundlagen identisch),
  \item $x$ das Alter bei Umstufung, $x_e$ das ursprüngliche
        Eintrittsalter,
  \item $u$ das in die Hilfsfunktion $b(\cdot)$ eingesetzte
        Umtarifierungs-Kostenargument in Vielfachen des Monatsbeitrags;
        es nimmt die Werte $u = 0$ (ohne Umtarifierungskosten) oder
        $u = \tilde{\alpha}$ (mit Umtarifierungskosten) an,
  \item $b^{\mathrm{a}}$ den bisher gezahlten monatlichen Beitrag,
        $b^{\mathrm{a}}_x$, $b^{\mathrm{n}}_x$ die Tarifbeiträge des Alters $x$ und
        $b^{\mathrm{n}}_{x_e}$ den Tarifbeitrag zum ursprünglichen
        Eintrittsalter,
  \item $\tilde{\alpha}$ den bei einer beitragserhöhenden Umstufung
        angesetzten Umtarifierungs-Kostensatz in Vielfachen des
        Monatsbeitrags; er ist auf $\tilde{\alpha} = 2$ festgelegt,
        während bei nicht beitragserhöhenden Umstufungen keine
        Umtarifierungskosten angesetzt werden ($\tilde{\alpha} = 0$).
\end{itemize}
bezeichnen. Die Hilfsfunktion $b(u)$ wird an den beiden
Kostenargumenten $u = 0$ und $u = \tilde{\alpha}$ ausgewertet; da
$b(u)$ je nach Vorzeichen monoton von $u$ abhängt, liefert
$b^{\mathrm{n}} = \max\{b(0);\, b(\tilde{\alpha})\}$ den für den Bestand
vorsichtigeren der beiden Randwerte. Die anschließende Kappung
$\hat{b}^{\mathrm{n}}$ begrenzt diesen nach oben auf den Tarifbeitrag $b^{\mathrm{n}}_x$
des erreichten Alters und nach unten auf den Neugeschäftsbeitrag
$b^{\mathrm{n}}_{x_e}$ zum ursprünglichen Eintrittsalter. Die Untergrenze
$b^{\mathrm{n}}_{x_e}$ stellt sicher, dass der neue Beitrag den
Neugeschäftsbeitrag zum ursprünglichen Eintrittsalter nicht
unterschreitet; ein hierdurch nicht verbrauchter Teil der angerechneten
Alterungsrückstellung wird der Rückstellung zur Beitragsermäßigung im
Alter zugeführt (\S\,13 Abs.\,1 KVAV).

Die Gleichartigkeit des Versicherungsschutzes richtet sich nach den
Leistungsbereichen des \S\,12 KVAV; ein Wechsel aus Tarifen mit
gleichartigem Versicherungsschutz erfolgt unter voller Anrechnung nach
Gleichung~\eqref{eq:umstufung}, bei nicht gleichartigem
Versicherungsschutz wird nur der auf die gemeinsamen Leistungsbereiche
entfallende Anteil der Alterungsrückstellung angerechnet. Entfällt bei
einer Umstufung ein Leistungsbereich (Teilstorno), wird die zugehörige
Alterungsrückstellung nicht angerechnet (\S\,12, \S\,13 Abs.\,2 KVAV).
Mittel aus dem gesetzlichen Beitragszuschlag
(Abschnitt~\ref{sec:zusatzrueckstellung}) und aus der Rückstellung nach
\S\,150 Abs.\,2 VAG verbleiben bei individuellen Tarifänderungen
innerhalb des Unternehmens ungekürzt dem Vertrag zugeordnet und werden
ab dem Alter 65 gemäß Gleichung~\eqref{eq:gutschrift}
beitragsmindernd eingesetzt.

\textbf{Wegfall oder Hinzunahme bestimmter Leistungen} ist im Tarif
VitalPlus nicht vorgesehen (kein Bausteintarif, für das Neugeschäft
geschlossen); Leistungsänderungen im Rahmen einer Beitragsanpassung
werden nach Abschnitt~\ref{sec:bap} behandelt. Ein gesonderter
Abschnitt entfällt daher.

\subsection{Beiträge für den Bestand bei Beitragsanpassungen}\label{sec:bap}

Bei einer Beitragsanpassung (Abschnitt~\ref{sec:af}) werden die
Beiträge des Bestandes gemäß \S\,11 KVAV nach den Grundsätzen der
Neuzugangskalkulation (Abschnitt~\ref{sec:beitraege}) mit den
aktualisierten Rechnungsgrundlagen neu berechnet; die vorhandene
Alterungsrückstellung wird vollständig prämienmindernd angerechnet.
Rechnerisch entspricht dies der Formel~\eqref{eq:umstufung} mit den
alten und neuen Rechnungsgrundlagen desselben Tarifs und
$\tilde{\alpha} = 0$ (in Prämien von Versicherten ab dem vollendeten
45.~Lebensjahr werden gemäß \S\,11 Abs.\,2 KVAV keine erneuten
einmaligen Abschlusskosten eingerechnet; das Unternehmen verzichtet
generell auf einmalige Kosten bei Beitragsanpassungen).
Risikozuschlagsfaktoren bleiben unverändert; absolute Selbstbehalte
werden nur im Rahmen des \S\,203 Abs.\,2 VVG angepasst.

\textbf{Limitierung:} Zur Begrenzung der Beitragssteigerung können
Mittel aus der Rückstellung für Beitragsrückerstattung eingesetzt
werden. Der limitierte Beitrag beträgt
\begin{equation}
b^{\mathrm{lim}} = \min\bigl\{\hat{b}^{\mathrm{n}};\;
b^{\mathrm{a}} + A_{\mathrm{lim}};\; b^{\mathrm{a}}\cdot(1+R_{\mathrm{lim}})\bigr\},
\label{eq:limit}
\end{equation}
wobei $\hat{b}^{\mathrm{n}}$ den formelmäßigen neuen Beitrag, $b^{\mathrm{a}}$ den
bisherigen Beitrag sowie $A_{\mathrm{lim}}$ (EUR) bzw.
$R_{\mathrm{lim}}$ (Prozentsatz) die vom Vorstand je
Beitragsanpassung festgelegten absoluten bzw. relativen
Limitierungsgrenzen bezeichnen. Die Beitragssenkung wird durch einen
Einmalbetrag in die Alterungsrückstellung gegenfinanziert:
\begin{equation}
E = 12\cdot\bigl(\hat{b}^{\mathrm{n}} - b^{\mathrm{lim}}\bigr)\cdot
\bigl(1-(\sigma+\Omega_x)\bigr)\cdot a_x,
\label{eq:einmalbetrag}
\end{equation}
wobei $x$ das erreichte Alter und $a_x$ den Rentenbarwert bezeichnen.

\subsection{Beitragsermäßigungen im Alter}\label{sec:beitragsermaessigung}

Mittel der Rückstellungen nach \S\,150 VAG sowie nicht verbrauchte
Teile des gesetzlichen Beitragszuschlags
(Abschnitt~\ref{sec:zusatzrueckstellung}) werden ab dem Alter 65 zur
Beitragsermäßigung eingesetzt; an über 80-jährige Versicherte werden
Zinserträge gemäß \S\,150 Abs.\,3 VAG außerhalb einer
Beitragsanpassung jährlich ausgekehrt. Ein Einmalbetrag $E$ mindert
den Monatsbeitrag um
\begin{equation}
b^{E}_x = \frac{E}{12\cdot\bigl(1-(\sigma+\Omega_x)\bigr)\cdot a_x},
\label{eq:gutschrift}
\end{equation}
wobei $x$ das Alter bei Anrechnung bezeichnet (keine Zillmerkosten bei
Gutschriften).
